\section{Behavioural Distances on LMCs}\label{sec:behdist}
We start with the definition of the type of system which we study in the remainder of the paper: labelled Markov chains.

\begin{definition}
	A labelled Markov chain (LMC) consists of the following data:
	\begin{itemize}
		\item a set of states \(X\);
		\item a (non-empty) set of labels \(L\);
		\item a (finitely branching) probabilistic transition function \(\tau \colon X \to \D_\mathbb{Q}(X)\); and
		\item a labelling function \(l \colon X \to L\)
	\end{itemize}
\end{definition}

\begin{example}
	Let \(X = \{x, x_1, x_2\}\) and \(L = \{a,b\}\). We represent the LMC \((X,L,\tau,l)\) with \(\tau(x) = \frac{1}{2} \ket{x_1} + \frac{1}{2} \ket{x_2}, \tau(x_1) = 1 \ket{x_1}, \tau(x_2) = 1 \ket{x_2}\) and \(l(x) = l(x_1) = a, l(x_2) = b\) as: %
	\begin{center}
		\begin{tikzpicture}[shorten >=0pt,node distance=2cm,
				on grid,auto,initial text=,every edge/.append style={bend angle=15}]
			\node[state] (x) {\slab{x}{a}} ;
			\node[state] (x1) [left=of x] {\slab{x_1}{a}} ;
			\node[state] (x2) [right=of x] {\slab{x_2}{b}} ;

			\path[->]
			(x) edge node[swap] {\(\frac{1}{2}\)} (x1)
			(x) edge node {\(\frac{1}{2}\)} (x2)

			;
		\end{tikzpicture}
	\end{center}
	We use the notation \slab{x}{a} for a state \(x \in X\) such that \(l(x) = a\). Further, any state with no outgoing edges is assumed to have a self-loop with probability \(1\).
\end{example}

We now recall a definition of the behavioural distance (henceforth written \(\bd\)) of states in an LMC as the least fixed point of a functional based on non-expansive maps.
This distinguishes two cases:
states having different labels should be maximally far apart, so they have distance \(1\);
the distance of states with the same label is then defined recursively, and can be seen as an optimisation problem.
Intuition for this problem is most often given in terms of its dual based on couplings under the Kantorovich-Rubinstein duality. The distance between distributions can in that setting be seen as the minimal cost of transporting one unit of mass from one distribution to the other, with the cost of transporting \(m\) probability mass between states at distance \(d\) being \(m \cdot d\).
In a more general form (as discussed in~\cite{villani2008optimal}) the distance below can be seen as the maximisation of profit, where we think of buying from one distribution and selling to the other, with the maps \(h\) representing buying/selling costs.

On small examples, these intuitions can often be used to determine the distance by examination; on larger systems a solver for linear programs is usually necessary.
For this, the definition below can be transformed into a rational linear program, as we use in the proof of \cref{lem:nstep}.
For further discussions of these distances, see for example~\cite{DBLP:journals/siglog/Breugel17, DBLP:conf/fossacs/RadyB23} %
\begin{definition}
	For \(X\) a set, and \(\PMet_X\) the set of pseudometric spaces on \(X\), we define \(\Gamma \colon \PMet_X \to \PMet_X\) by
	\begin{align*}
		\Gamma(d)(x,y) = \begin{cases}
			                 1,                                                                           & \text{ if } l(x) \neq l(y), \\
			                 \sup_{h \colon (X,d) \to ([0,1],d_e)} \tau(x) \cdot h - \tau(y) \cdot h, & \text{ o.w.}
		                 \end{cases}
	\end{align*}
	Then we define \(\bd := \lfp(\Gamma)\).
\end{definition}
Note that the least fixed point exists, because \(\PMet_X\) is a complete lattice, and \(\Gamma\) preserves the pointwise order on \(\PMet_X\).

\begin{example}\label{ex:LMC}
	Consider the following LMC:
	\begin{center}
		\begin{tikzpicture}[shorten >=0pt,node distance=2cm,
				on grid,auto,initial text=,every edge/.append style={bend angle=15}]
			\node[state] (x) {\slab{x}{a}} ;
			\node[state] (x1) [left=of x] {\slab{x_1}{a}} ;
			\node[state] (x2) [right=of x] {\slab{x_2}{b}} ;
			\node[state] (y) [right=7cm of x] {\slab{y}{a}} ;
			\node[state] (y1) [left=of y] {\slab{y_1}{a}} ;
			\node[state] (y2) [right=of y] {\slab{y_2}{b}} ;


			\path[->]
			(x) edge node[swap] {\(\frac{1}{2}\)} (x1)
			(x) edge node {\(\frac{1}{2}\)} (x2)
			(y) edge node[swap] {\(\frac{2}{5}\)} (y1)
			(y) edge node {\(\frac{3}{5}\)} (y2)
			;
		\end{tikzpicture}
	\end{center}
	Note that \(\bd(x_1,y_1) = \bd(x_2,y_2) = 0\) and \(\bd(x_1,y_2) = \bd(x_2,y_1) = 1\), which are the values given by \(\Gamma(\bot)\). The value \(\bd(x,y)\) is then \(\Gamma^2(\bot)(x,y)\) for which it can be shown that the supremum is achieved by the map \(h_0(z) = \textbf{if } z \in \{x_1,y_1\} \textbf{ then } 1 \textbf{ else } 0\) so that:
	\begin{align*}
		\Gamma^2(\bot)(x,y) & = \sup_{h \colon (X,\Gamma(\bot)) \to ([0,1],d_e)} \tau(x) \cdot h - \tau(y) \cdot h                                         \\
		                    & = \tau(x) \cdot h_0 - \tau(y) \cdot h_0                                                                                      \\
		                    & = \left(\frac{1}{2} \cdot 1 + \frac{1}{2} \cdot 0\right) - \left(\frac{2}{5} \cdot 1 + \frac{3}{5} \cdot 0\right) = \frac{1}{10}
	\end{align*}
	In terms of transportation costs, this \(\frac{1}{10}\) can be seen as the cost of transporting the \(\frac{1}{10}\) of mass from \(x_1\) to \(y_2\) which can not be transported to \(y_1\) due to its lack of capacity.
	Intuition based on profits is harder to give in this case, due to the asymmetry of our definition. It is easier to obtain the \(h\) in this example by seeing it as grouping the states into ``equivalence'' classes, spaced as far apart as possible while respecting the distance on \(X\) (in this case \(\Gamma(\bot)\)).
\end{example}

It will be important for the correspondence results of later sections that the behavioural distance can be obtained as a countable supremum, namely the supremum over all finite applications of \(\Gamma\) to the constant zero distance.
A similar result for LMCs with non-determinism is shown already in~\cite[Sec.~3]{DBLP:conf/lics/DesharnaisJGP02}.
It can also be proved using the Kleene fixpoint theorem, or
\(\omega\)-(co)continuity of \(\Gamma\) as shown in~\cite{DBLP:journals/ipl/Breugel12}.
\begin{proposition}\label{prop:kleene}
	For any LMC \((X,L,\tau,l)\) and \(x,y \in X\), we have
	\begin{equation*}
		\bd(x,y) = \sup_{i < \omega} \Gamma^i(\bot)(x,y)
	\end{equation*}
\end{proposition}

\section{Proof System}\label{sec:proofsys}

In this section, we define our derivation system for lower bounds on behavioural distances between states of an LMC.
The conclusion of the rules are of the form \(x \apart{\varepsilon} y\) which, as our soundness result will show, implies that \(\bd(x,y) \geq \varepsilon\), i.e., the behavioural distance between \(x\) and \(y\) is at least \(\varepsilon\).
The definition of \(\bd\) suggests two rules, one for each case. The label case straightforwardly yields the rule
\begin{equation*}
	\AxiomC{\(l(x) \neq l(y)\)}
	\RightLabel{(lab)}
	\UnaryInfC{\(x \apart{1} y\)}
	\DisplayProof
\end{equation*}
In the supremum case, \(\bd(x,y)\) can be bounded from below by \(\tau(x) \cdot h - \tau(y) \cdot h\) for any non-expansive map \(h\) by definition.
However, it is not immediately clear that this can be done inductively, as we can not assume \(\bd\) to be known, and thus cannot use it to choose a non-expansive \(h\).
Fortunately, as long as the system is sound with respect to the behavioural distance, it suffices to have a map \(h\) for which a kind of \emph{pairwise non-expansiveness} holds: for any \(x',y'\) we have \(|h(x') - h(y')| \leq \varepsilon\) for some \(\varepsilon\) such that we have proved \(x' \apart{\varepsilon} y'\). Soundness then implies that \(|h(x') - h(y')| \leq \varepsilon \leq \bd(x',y')\) for all \(x',y'\), which is exactly non-expansiveness of \(h\) with respect to \(\bd\).

Now, in proofs, we could allow arbitrary recursive proofs and require the choice of a pairwise non-expansive map to correctly apply the rule. Alternatively, we can choose to allow arbitrary maps.
We are then required to prove that for all \(x',y' \in X\), \(|h(x') - h(y')|\) is a lower bound on the behavioural distance.
We can see the corresponding proof obligations \(x' \apart{|h(x') - h(y')|} y'\) as those generated by a chosen map \(h\). The first option fits with a forward reasoning approach to constructing a proof; we prove some bounds and try to find a fitting \(h\). The second is a backward approach; if we wish to show a bound \(x \apart{\varepsilon} y\), we must supply an \(h\) and recursively prove its validity.

We choose the latter approach, primarily because it makes the proof obligations clearer, and we will be able to see when a choice of map is not valid. Using the earlier form, a chosen \(h\) may be invalid because we have not proved strong enough bounds, or because it is simply not non-expansive with respect to \(\bd\).
Such a rule can be written as follows:
\begin{equation*}
	\AxiomC{\(h \colon X \to [0,1]\)}
	\AxiomC{\(\forall x', y' \in X \ldotp x' \apart{|h(x') - h(y')|} y'\)}
	\AxiomC{\(\tau(x) \cdot h - \tau(y) \cdot h \geq \varepsilon\)}
	\TrinaryInfC{\(x \apart{\varepsilon} y\)}
	\DisplayProof
\end{equation*}
However, in this form, the rule does not give a usable derivation system.
Namely, in case an LMC has an infinite state space, there will be infinitely many recursive proof obligations.
We will ensure that this is always finite, and even reduce the work required to construct proofs beyond this.
Further, the map \(h\) is so far \emph{real}-valued. This poses a problem both for our approximate completeness result, in which we need to be able to compute these maps, and the construction from proofs to modal formulas whose interpretation will be rational-valued.
Summarising, to remedy these issues, we adapt the above rule in the following ways:
\begin{itemize}
	\item we show that \(h\) need only be defined on a finite subset of the state space thereby generating only finitely many proof obligations;
	\item we restrict the codomain of \(h\) to rationals;
	\item we reduce the number of recursive proof obligations further by not requiring proofs for those bounds which follow from reflexivity and symmetry.
\end{itemize}

To make our proof system and its presentation more pleasant, we include three more rules inspired by those from quantitative equational theories~\cite{DBLP:conf/lics/MardarePP16}. Namely: a \emph{zero} (reflexivity) rule; a \emph{symmetry} rule; and a \emph{weakening} rule.
Together, this brings us to the following rules:
\begin{definition}\label{def:proofrules}
	Let \((X,L,\tau,l)\) be an LMC, \(x,y \in X\), and \(\varepsilon \in \UIQ\). Further, we define \(\supps_{x,y} := \{ s \in X \mid \tau(x)(s) \neq \tau(y)(s) \}\) and \(\mu \cdot_\supps h := \sum_{s \in \supps} \mu(s) \cdot h(s)\). We may drop the subscripts \(x,y\) and \(\supps\) whenever clear from the context.

	Then, we define the following derivation rules:
	\begin{equation*}
		\AxiomC{\vphantom{\(y \apart{\varepsilon} x\)}}
		\RightLabel{(zero)}
		\UnaryInfC{\(x \apart{0} y\)}
		\DisplayProof
		\qquad
		\AxiomC{\(y \apart{\varepsilon} x\)}
		\RightLabel{(symm)}
		\UnaryInfC{\(x \apart{\varepsilon} y\)}
		\DisplayProof
		\qquad
		\AxiomC{\(x \apart{\varepsilon'} y\)}
		\AxiomC{\(\varepsilon \leq \varepsilon'\)}
		\RightLabel{(weak)}
		\BinaryInfC{\(x \apart{\varepsilon} y\)}
		\DisplayProof
		\qquad
		\AxiomC{\(l(x) \neq l(y)\)}
		\RightLabel{(lab)}
		\UnaryInfC{\(x \apart{1} y\)}
		\DisplayProof
	\end{equation*}
	\begin{equation*}
		\AxiomC{\(h \colon \supps \to \UIQ\)}
		\noLine
		\UnaryInfC{\(\forall x', y' \in \supps \ldotp h(x') > h(y') \implies x' \apart{h(x') - h(y')} y'\)}
		\AxiomC{\(\tau(x) \cdot_\supps h - \tau(y) \cdot_\supps h \geq \varepsilon\)}
		\RightLabel{(exp)}
		\BinaryInfC{\(x \apart{\varepsilon} y\)}
		\DisplayProof
	\end{equation*}
	We write \(\PTs{X}\) for the smallest set which contains all instances of the \emph{(zero)} and \emph{(lab)} rules for \(x,y \in X\), and is closed under applications of all instances of \emph{(symm)}, \emph{(weak)}, and \emph{(exp)} for any \(x,y \in X\) and \(\varepsilon \in \UIQ\).
\end{definition}
Note that we write conditions as premises in the rules rather than as separate side-conditions for convenience. As suggested by the definition of \(\PTs{X}\), this makes both the \emph{(zero)} and \emph{(lab)} rules axioms.
Further, strictly speaking, \emph{(exp)} is a family of rules, indexed by the maps \(h\).
The set \(\PTs{X}\) can be thought of as the set of all proof trees which can be built from the given rules.
As is usual, we will write \(\vdash x \apart{\varepsilon} y\) to mean that the given judgment is provable, i.e., there is a proof tree in \(\PTs{X}\) with the given judgment at the root.
\begin{remark}
	Note that the restriction to \(\supps\) in \emph{(exp)} means proofs, which are finite depth by definition, will also be finite breadth even when the LMC under consideration has an infinite state space.
	This is because \(\{ s \in X \mid \tau(x)(s) \neq \tau(y)(s) \}\) is a subset of \(\supp(\tau(x)) \cup \supp(\tau(y))\), which is finite by assumption.
	We will see later that this allows us to provide witnesses as both finite proof trees and finite modal formulas.
	This improves on the earlier work of~\cite{DBLP:conf/fossacs/RadyB23} which restricts to finite state spaces.
	We illustrate this improvement in Example~25 in the original paper (p.~25:14), once we have shown soundness and completeness of the proof system, and its correspondence with modal formulas.
\end{remark}

\begin{example}\label{ex:simpleproof}
	We continue with the LMC from \cref{ex:LMC} and show how we can prove the distance between \(x\) and \(y\) shown there as a lower bound. Using the \emph{(lab)} rule, the bounds \(u \apart{1} v\) can be proved for \(u,v \in \{x_1,y_1\}\) and \(v \in \{x_2,y_2\}\).
	This allows us to define \(h_0 \colon \supps \to \UIQ\) as before by \(h_0(z) = \textbf{if } z \in \{x_1,y_1\} \textbf{ then } 1 \textbf{ else } 0\) for which \(\tau(x) \cdot h_0 - \tau(y) \cdot h_0 = \frac{1}{10}\) so that we can prove
	%
	\begin{prooftree}%
		\AxiomC{}
		\UnaryInfC{\(x_1 \apart{1} x_2\)}
		\AxiomC{}
		\UnaryInfC{\(x_1 \apart{1} y_2\)}
		\AxiomC{}
		\UnaryInfC{\(y_1 \apart{1} x_2\)}
		\AxiomC{}
		\UnaryInfC{\(y_1 \apart{1} y_2\)}
		\QuaternaryInfC{\(x \apart{\frac{1}{10}} y\)}
	\end{prooftree}
\end{example}

\begin{example}\label{ex:noform}
	Our second example serves to illustrate a limitation of our proof system, namely that the behavioural distance of states will not always be exactly provable in our system. It would only be provable if we allowed infinite depth proof trees.
	The LMC we consider is the same as the one in~\cite[Thm.~17]{DBLP:conf/fossacs/RadyB23}, which shows that there is an LMC containing states for which it is not possible to give a single formula ``explaining'' their distance.
	\begin{center}
		\begin{tikzpicture}[shorten >=0pt,node distance=2cm,
				on grid,auto,initial text=,every edge/.append style={bend angle=15}]
			\node[state] (x) {\slab{x}{a}} ;
			\node[state] (x1) [left=of x] {\slab{x_1}{b}} ;
			\node[state] (y) [right=5cm of x] {\slab{y}{a}} ;


			\path[->]
			(x) edge node[swap] {\(\frac{1}{2}\)} (x1)
			(x) edge[loop right] node[swap] {\(\frac{1}{2}\)} (x)
			;
		\end{tikzpicture}
	\end{center}
	As is discussed in \emph{op.\ cit.}, the distance \(\bd(x,y) = 1\) is reached only in the limit, not by any \(\Gamma^i(\bot)\) and thus not by any single tree.
	Proving the bound given by \(\Gamma^i(\bot)\) (for \(i > 0\)) can be done using \(i - 1\) applications of the \emph{(exp)} rule together with two applications each of the \emph{(lab)} and \emph{(zero)} rules.
	For example, once we have proved \(x_1 \apart{1} u\) for \(u \in \{x,y\}\), we can take the map \(h_0 \colon x,y \mapsto 0, x_1 \mapsto 1\) for which \(\tau(x) \cdot h_0 - \tau(y) \cdot h_0 = \frac{1}{2}\) and prove:
	%
	\begin{prooftree}
		\AxiomC{}
		\UnaryInfC{\(x_1 \apart{1} x\)}
		\AxiomC{}
		\UnaryInfC{\(x_1 \apart{1} y\)}
		\BinaryInfC{\(x \apart{\frac{1}{2}} y\)}
	\end{prooftree}
	This step (plus an application of \emph{(symm)}) allows the next application of \emph{(exp)} with a non-expansive \(h_0\) mapping \(x\) to \(\frac{1}{2}\), yielding a bound of \(\frac{3}{4}\). Continuing to increase the value of \(h_0(x)\) in this way, we approach \(\bd(x,y)\) from below.
\end{example}

\subsection{Soundness}\label{sec:sound}
We now move on to showing soundness of the system, i.e., that if we can prove \(x \apart{\varepsilon} y\), then \(\bd(x,y) \geq \varepsilon\).
The \emph{(zero)} rule is sound as our pseudometrics are valued in \([0,1]\) and thus 0 is always a sound lower bound.
Similarly, the behavioural distance is symmetric, so that the order of states does not change a lower bound's validity.
Our weakening rule is sound by transitivity of \(\leq\).
Soundness of the label rule follows from the definition of \(\Gamma\). This is similar for the expectation rule, but this requires some more work. The discussion at the beginning of this section gives some intuition.

Due to our restriction of the domain of the map in the \emph{(exp)} rule, we will require the following lemma in the soundness proof:
\begin{restatable}{lemma}{supsupp}\label{lem:supsupp}
	For \((X,L,\tau,l)\) an LMC, \(d \colon X \times X \to [0,1]\) a pseudometric and \(x,y \in X\):
	\[
		\sup_{h \colon (X,d) \to ([0,1],d_e)} \tau(x) \cdot h - \tau(y) \cdot h = \sup_{h \colon (\supps,d|_\supps) \to ([0,1],d_e)} \tau(x) \cdot_\supps h - \tau(y) \cdot_\supps h
	\]
	where \(d|_\supps = d \circ (\iota_\supps \times \iota_\supps)\) with \(\iota_\supps \colon \supps \hookrightarrow X\) the inclusion map.
\end{restatable}
\begin{proofappendix}{supsupp}
	We prove two inequalities:
	\subparagraph*{\(\leq\):}
	This holds because any \(h \colon (X,d) \to ([0,1],d_e)\) restricts to a map \(h|_\supps = h \circ \iota_\supps \colon (\supps,d|_\supps) \to ([0,1],d_e)\), and we can show that
	\[
		\tau(x) \cdot h - \tau(y) \cdot h = \tau(x) \cdot_\supps h|_\supps - \tau(y) \cdot_\supps h|_\supps
	\]
	\subparagraph*{\(\geq\):}
	For this direction, we show that for any \(h \colon (\supps, d|_\supps) \to ([0,1],d_e)\), there is an \(h' \colon (X,d) \to ([0,1],d_e)\) such that
	\[
		\tau(x) \cdot h' - \tau(y) \cdot h' = \tau(x) \cdot_\supps h - \tau(y) \cdot_\supps h
	\]
	We use an existing construction of extensions of non-expansive maps, to extend \(h\) along the inclusion \(\iota_\supps \colon \supps \hookrightarrow X\). Namely, we define \(h'(x) := \inf_{z \in \supps} h(z) \oplus d(x,z)\), where \(\oplus\) is truncated addition on the unit interval.
	This is an extension in the sense that \(h' \circ \iota_\supps = h\), so that the above equality indeed holds.
\end{proofappendix}

\begin{remark}\label{rem:hfromLP}
	We can further restrict the space of possible maps \(h\) in the above, by seeing the supremum as the optimal solution of a linear program as follows.
	We encode functions \(h \colon \supps \to [0,1]\) as (finite) vectors \(\vec{h} \in [0,1]^{|\supps|}\), writing \(\vec{h}_x\) for the value of \(\vec{h}\) at the position indexed by \(x \in \supps\).
	Then, each inequality \(|h(x) - h(y)| \leq d(x,y)\) can be expressed by \(\vec{a} \cdot \vec{h} \leq d(x,y)\) and \(\vec{a}' \cdot \vec{h} \leq d(x,y)\) with \(\vec{a}_x = 1, \vec{a}_y = -1, \vec{a}'_x = -1, \vec{a}'_y = 1\) (and all other entries zero). We can enforce \(0 \leq \vec{h}_x \leq 1\) similarly.
	We are thus interested in the problem of maximising \(\tau(x) \cdot h - \tau(y) \cdot h\) (a linear expression) subject to the constraints expressed by \(\mathbf{A} \cdot \vec{h} \leq \vec{b}\) for an integer matrix \(\mathbf{A}\) and vector \(\vec{b}\).

	The feasible region of this problem is a polytope, which we call \(\Hpoly\). It is convex and closed due to the shape of the constraints, and bounded due to the restriction to the unit interval. It is furthermore non-empty as any constant map valued in \([0,1]\) lies within it.
	It is known (see, e.g.,~\cite{trustrum1971linear}) that for such a feasible region, the optimal value is achieved in the vertices of the polytope, which we denote by \(V(\Hpoly)\). Then we can write
	\begin{equation*}
		\sup_{h \colon (X,d) \to ([0,1],d_e)} \tau(x) \cdot h - \tau(y) \cdot h = \max_{h \in V(\Hpoly)} \tau(x) \cdot_\supps h - \tau(y) \cdot_\supps h
	\end{equation*}
	The supremum becomes a maximum because finitely many linear inequalities define a polytope with finitely many vertices.

	Obtaining a map \(h\) in which the supremum is achieved can thus be done using an algorithm such as simplex. We will apply this recursively for \(d = \Gamma^i(\bot)\) (the finite approximants of \(\bd\)), to construct proof trees in our proof of approximate completeness in the next section.
\end{remark}


We are now able to show (by structural induction) that any proof in our system yields a lower bound on \(\bd\).
\begin{restatable}[Soundness]{theorem}{sound}\label{lem:sound}
	For any LMC \((X,L,\tau,l)\), any proof tree built from the rules of \cref{def:proofrules}, any \(\varepsilon \in \UIQ\), and any \(x, y \in X\), if the proof tree has \(x \apart{\varepsilon} y\) at the root, then \(\bd(x,y) \geq \varepsilon\).
\end{restatable}
\begin{proofappendix}{sound}
	We proceed by induction on the structure of the proof tree.

	\textbf{Case \emph{(zero)}:} By its definition, \(\bd\) takes values in \([0,1]\), so that \(\bd(x,y) \geq 0\) holds.

	\textbf{Case \emph{(label)}:} We have a proof tree
	\[
		\AxiomC{\(l(x) \neq l(y)\)}
		\UnaryInfC{\(x \apart{1} y\)}
		\DisplayProof
	\]
	By definition of \(\Gamma\), we must have \(\bd(x,y) = 1\), so that indeed \(\bd(x,y) \geq 1\).

	\textbf{Case \emph{(symm)}:} We have a proof tree
	\[
		\AxiomC{\(y \apart{\varepsilon} x\)}
		\UnaryInfC{\(x \apart{\varepsilon} y\)}
		\DisplayProof
	\]
	By induction, we have \(\bd(y,x) \geq \varepsilon\), but \(\bd\) is symmetric, so that also \(\bd(x,y) \geq \varepsilon\).

	\textbf{Case \emph{(weak)}:} We have a proof tree
	\[
		\AxiomC{\(x \apart{\varepsilon'} y\)}
		\AxiomC{\(\varepsilon \leq \varepsilon'\)}
		\BinaryInfC{\(x \apart{\varepsilon} y\)}
		\DisplayProof
	\]
	By induction, we have \(\bd(x,y) \geq \varepsilon' \geq \varepsilon\).

	\textbf{Case \emph{(exp)}:} We have a proof tree
	\begin{equation*}
		\expr
	\end{equation*}
	By induction, we have for all \(x',y' \in \supps\) with \(h(x') > h(y')\) that \(\bd(x',y') \geq |h(x') - h(y')|\). For \(x',y' \in \supps\) with \(h(x') < h(y')\), we have \(\bd(x',y') = \bd(y',x') \geq |h(y') - h(x')| = |h(x') - h(y')|\). For the remaining pairs, \(|h(x') - h(y')| = 0 \leq \bd(x',y')\). In other words, \(h \colon \supps \to \UIQ\) is a non-expansive map \(h \colon (\supps,\bd) \to (\UIQ,d_e)\).
	As \(\bd\) is defined as the least fixed point of \(\Gamma\), we have
	\begin{align*}
		\bd(x,y) & = \Gamma(\bd)(x,y)                                                                                                              \\
		         & = \begin{cases}
			             1,                                                                             & \text{ if } l(x) \neq l(y), \\
			             \sup_{h \colon (X,\bd) \to ([0,1],d_e)} \tau(x) \cdot h - \tau(y) \cdot h, & \text{ o.w.}
		             \end{cases}
	\end{align*}
	In case \(l(x) \neq l(y)\), we have \(\bd(x,y) = 1 \geq \varepsilon\).

	In the remaining case, we have
	\begin{align*}
		\bd(x,y) & = \sup_{k \colon (X,\bd) \to ([0,1],d_e)} \tau(x) \cdot k - \tau(y) \cdot k                          \\
		         & = \sup_{h \colon (\supps,d|_\supps) \to ([0,1],d_e)} \tau(x) \cdot_\supps h - \tau(y) \cdot_\supps h \\
		         & \geq \tau(x) \cdot h - \tau(y) \cdot h \geq \varepsilon
	\end{align*}
	where the second equality is shown in \cref{lem:supsupp} and the first inequality holds because \(h\) is one of the non-expansive maps ranged over in the \(\sup\).
	This covers all cases, so soundness follows by induction.
\end{proofappendix}

\subsection{Approximate Completeness}\label{sec:appcomp}
The rest of this section is dedicated to proving the approximate completeness of the system.
This will show that we can prove lower bounds arbitrarily close to the ``true'' value given by \(\bd\).
Our proof relies on the fact that we can get arbitrarily close to \(\bd\) with its finite approximants \(\Gamma^i(\bot)\), and the following lemma, that shows how we can construct proofs exhibiting these finite approximants as lower bounds on the behavioural distance.

\begin{lemma}\label{lem:nstep}
	For any \(i \in \mathbb{N}\) and \(x,y \in X\), we have \(\vdash x \apart{\Gamma^i(\bot)(x,y)} y\).
\end{lemma}
\begin{proof}
	By induction on \(i\), where we strengthen the induction by showing that \(\Gamma^i(\bot)(x,y)\) is always rational.
	For the base case, we have \(\Gamma^0(\bot)(x,y) = \bot(x,y) = 0\) which is rational, and we can prove \(x \apart{0} y\) using the \emph{(zero)} rule.

	Now let \(i \in \mathbb{N}\), and suppose for any \(x,y \in X\), that we can prove \(x \apart{\Gamma^i(\bot)(x,y)} y\) and that \(\Gamma^i(\bot)(x,y)\) is rational. We have
	\begin{align*}
		\Gamma^{i+1}(\bot)(x,y) & = \Gamma(\Gamma^i(\bot))(x,y)                                                                                                     \\
		                        & = \begin{cases}
			                            1,                                                                                        & \text{ if } l(x) \neq l(y), \\
			                            \sup_{h \colon (X,\Gamma^i(\bot)) \to ([0,1],d_e)} \tau(x) \cdot h - \tau(y) \cdot h, & \text{ o.w.}
		                            \end{cases}
	\end{align*}
	If \(l(x) \neq l(y)\), we can prove \(x \apart{1} y\) using the \emph{(lab)} rule.
	Otherwise, by \cref{lem:supsupp}, we have
	\begin{align*}
		\sup_{h \colon (X,\Gamma^i(\bot)) \to ([0,1],d_e)} \tau(x) \cdot h - \tau(y) \cdot h = \sup_{h \colon (\supps,\Gamma^i(\bot)|_\supps) \to ([0,1],d_e)} \tau(x) \cdot_\supps h - \tau(y) \cdot_\supps h
	\end{align*}
	As explained in \cref{rem:hfromLP}, we can find an optimal \(h_0\) by solving a linear program.
	This will in particular be a map \(h_0 \colon (\supps,\Gamma^i(\bot)|_\supps) \to (\UIQ,d_e)\) (note the restriction to rationals) because all the coefficients in the problem are rational, by induction. We thus have
	\[
		\sup_{h \colon (\supps,\Gamma^i(\bot)|_\supps) \to ([0,1],d_e)} \tau(x) \cdot_\supps h - \tau(y) \cdot_\supps h = \tau(x) \cdot_\supps h_0 - \tau(y) \cdot_\supps h_0
	\]
	which is rational, because the transition probabilities given by \(\tau(x)\) and \(\tau(y)\) are also rational.
	We can now construct the following proof, in which recursive proofs are given by induction, and the above discussion allows us to choose \(\varepsilon := \Gamma^{i+1}(\bot)(x,y)\):
	\begin{equation*}
		\AxiomC{\(h_0 \colon \supps \to \UIQ\)}
		\noLine
		\UnaryInfC{\(\forall x', y' \in \supps \ldotp h_0(x') > h_0(y') \implies x' \apart{h_0(x') - h_0(y')} y'\)}
		\AxiomC{\(\tau(x) \cdot h_0 - \tau(y) \cdot h_0 \geq \varepsilon\)}
		\BinaryInfC{\(x \apart{\varepsilon} y\)}
		\DisplayProof
	\end{equation*}\qedhere
\end{proof}
We see that completeness in this sense only requires the use of the \emph{(zero)}, \emph{(lab)}, and \emph{(exp)} rules. In fact, it can be shown for only the latter two, with the \emph{(zero)} rule being essentially an instance of the \emph{(exp)} rule where the map \(h\) is taken to be constant. The approximate completeness of the system is now a simple consequence of the above lemma and \cref{prop:kleene}, with no further applications of the rules needed.
\begin{restatable}[Approximate Completeness]{theorem}{approxcomp}\label{thm:approxcomp}
	For any (real) \(\delta > 0\), and any \(x,y \in X\) there is a proof tree with \(x \apart{\varepsilon} y\) at the root, so that \(0 \leq \bd(x,y) - \varepsilon < \delta\).
\end{restatable}
\begin{proofappendix}{approxcomp}
	Let \(\delta > 0\). By \cref{prop:kleene}, we have \(\bd(x,y) = \sup_{i < \omega} \Gamma^i(\bot)(x,y)\), so there exists \(i \in \mathbb{N}\) such that
	\[
		0 \leq \bd(x,y) - \Gamma^i(\bot)(x,y) < \delta
	\]
	\Cref{lem:nstep} exactly gives us a proof of \(x \apart{\Gamma^i(\bot)(x,y)} y\), and we are done.
\end{proofappendix}
